\chapter{概述}\label{ch:1-overview}

\section{简介}\label{sec:1-1}

在现代微控制器（MCU）应用中，人们对灵活性、安全性和可靠性的要求日益提高。越来越多的微控制器应用要求多个互不信任的来源的应用程序一起运行，也可能同时要求实时可靠部分\footnote{如电机控制等}、性能优先部分\footnote{如网络协议栈、高级语言虚拟机等}和某些对信息安全有要求的应用\footnote{如加解密算法等}一起运行，而且互相之间不能干扰。在某些应用程序宕机时，要求它们能够单独重启，并且在此过程中还要求保证那些直接控制物理机电输出的应用程序不受干扰。

由于传统实时操作系统（Real-Time Operating
System，RTOS）的内核和应用程序通常静态链接在一起，而且互相之间没有内存保护隔离，因此达不到应用间信息隔离的要求；一旦其中一个应用程序崩溃而破坏内核数据，其他所有应用程序必然同时崩溃。一部分RTOS号称提供了内存保护，但是其具体实现机理表明其内存保护仅仅能够提供一定程度的应用间可靠性而无法提供应用间安全性，不适合现代微控制器应用的信息安全要求。此外，传统RTOS的内核服务\footnote{如内核定时器等}是在高优先级应用和低优先级应用之间共用的，一旦低优先级应用调用大量的内核服务，势必影响高优先级应用内核服务的响应，从而造成时间干扰问题\protect\hyperref[cite:1-1]{{[}1{]}}。

RVM（Real-time Virtual machine
Monitor）就是针对上述问题而提出的、面向高性能MCU的、基于RME微内核的用户态库。它是运行在RME操作系统上的一个宿主型虚拟机监视器，使得在微控制器平台上运行多个虚拟机成为可能，也同时满足用户运行原生RME应用程序的需求。这样，实时和非实时应用一起运行的需求，在实时应用和非实时应用之间不会产生时间干扰。它也能够使多来源应用程序在同一芯片上运行而不发生信息安全问题。

由于RVM强制隔离了系统中安全的部分和不安全的部分，这两部分就可以用不同的标准分开开发和认证。这方便了调试，也节省了开发和认证成本，尤其是当使用到那些较为复杂、认证程度较低和实时性较差的的软件包如用户图形界面（Graphical
User
Interface，GUI）和高级语言虚拟机（Python、Java虚拟机）等的时候。此外，在老平台上已得到认证的微控制器应用程序可以作为一个虚拟机直接运行在RVM上，无需重新认证；这进一步节省了成本。

本手册从用户的角度提供了RVM的功能描述。关于各个架构的具体使用，请参看各个架构相应的手册。在本手册中，我们先简要回顾关于虚拟化的若干概念，然后分章节介绍RVM的特性和API。

\subsection{微控制器的特点}\label{subsec:1-1-1}

在实际应用中，MCU系统占据了实时系统中的相当部分。它常常有如下几个特点：

\subsubsection{深度嵌入}\label{subsubsec:1-1-1-1}

微控制器的应用场合往往和设备本身融为一体，用户通常无法直接感知到此类系统的存在。微控制器的应用程序\footnote{通常称为``固件''}一旦写入，往往在整个生命周期中不再更改，或者很少更改。即便有升级的需求，往往也由厂家进行升级。

\subsubsection{高度可靠}\label{subsubsec:1-1-1-2}

微控制器应用程序对于可靠性有近乎偏执的追求，任何系统故障都会引起相对严重的后果。从小型电压力锅到大型起重机，这些系统都不允许发生严重错误，否则会造成重大财产损失或人身伤亡。此外，对于应用的实时性也有非常高的要求，它们通常要么要求整个系统都是硬实时系统，或者系统中直接控制物理系统的那部分为硬实时系统。

\subsubsection{资源有限}\label{subsubsec:1-1-1-3}

由于微控制器本身资源不多，因此高效地利用它们就非常重要了。微控制器所使用的软件无论是在编写上和编译上总体都为空间复杂度优化，只有小部分对性能和功能至关重要的程序使用时间复杂度优化。

\subsubsection{静态分配和链接}\label{subsubsec:1-1-1-4}

与微处理器不同，由于深度嵌入的原因，微控制器应用中的绝大部分资源，包括内存和设备在内，都是在系统上电时创建并分配的。通常而言，在微控制器内使用过多的动态特性是不明智的，因为动态特性不仅会对系统的实时响应造成压力，另一方面也会增加功能成功执行的不确定性。如果某些至关重要的功能和某个普通功能都依赖于动态内存分配，那么当普通功能耗尽内存时，可能导致重要功能的内存分配失败。基于同样的或类似的理由，在微控制器中绝大多数代码都是静态链接的，一般不使用动态链接。

\subsubsection{仅物理地址空间}\label{subsubsec:1-1-1-5}

与微处理器不同，微控制器通常都不具备内存管理单元（Memory Management
Unit，MMU），因此无法实现物理地址到虚拟地址的转换。但是，RVM也支持使用内存保护单元（Memory
Protection
Unit，MPU），并且大量的中高端微控制器都具备MPU。与MMU不同，MPU往往仅支持保护一段或几段内存范围，有些还有很复杂的地址对齐限制，因此内存管理功能在微控制器上往往是有很大局限性的。

\subsubsection{单核处理器}\label{subsubsec:1-1-1-6}

如今，绝大多数的微处理器都是多核设计。然而，与微处理器不同，微控制器通常均为单核设计，因此无需考虑很多竞争冒险问题。这可以进一步简化微控制器用户态库的设计。考虑到的确有少部分微控制器采取多核设计或者甚至不对称多核设计，RVM也支持多核运行环境。当然，这首先需要底层的RME被正确配置为支持多核。

考虑到以上几点，在设计RVM时，所有的内核对象都会在系统启动时创建完全，并且在此时，系统中可用内核对象的上限也就决定了。各个虚拟机则被固化于微控制器的片上非易失性存储器中。

\subsection{软件版权与许可证}\label{subsec:1-1-2}

综合考虑到微控制器应用、深度嵌入式应用和传统应用对开源系统的不同要求，RVM微控制器库所采用的许可证为Unlicensed，但是对一些特殊情况\footnote{比如安防器材、军工系统、航空航天装备、电力系统和医疗器材等}使用特殊的规定。这些特殊规定是就事论事的，对于每一种可能情况的具体条款都会有不同。

\subsection{主要参考系统}\label{subsec:1-1-3}

RTOS通用服务主要参考了如下系统的实现：

\begin{quote}
RMProkaron（@EDI）

FreeRTOS（@Real-Time Engineering LTD）

Contiki（@The Contiki Community）
\end{quote}

\needspace{8\baselineskip}
Contiki虚拟机监视器的设计参考了如下系统的实现：

\begin{quote}
Nova（@TU Dresden）

XEN（@Cambridge University）

KVM（@The Linux Foundation）
\end{quote}

其他各章的参考文献和参考资料在该章列出。

\section{前言}\label{sec:1-2}

作为一个虚拟机监视器，RVM依赖于底层RME操作系统提供的内核服务来维持运行。RVM采用准虚拟化技术，最大限度节省虚拟化开销，典型情况下仅增加约1\%的CPU开销。

RVM完全支持低功耗技术，可以实现虚拟化系统的各个分区处于无节拍工作，最大限度减少电力需求。

对于RME原生进程应用，其响应时间和FreeRTOS一档小型RTOS的响应时间相若，完全能够满足那些对实时性要求最苛刻的应用。

对于RVM宿主型虚拟机\footnote{虚拟机分类见下文所述}，中断响应时间和线程切换时间将会增加到原小型RTOS的约4倍。值得注意的是，RVM本身并不提供任何POSIX支持；POSIX支持是依赖于所移植的操作系统本身的。幸运的是，很多RTOS都具备POSIX支持，因此直接虚拟化它们就可以使用其POSIX功能了。

RME和RVM在全功能配置下共占用最小16kB内存\footnote{此时可支持4个虚拟机，若分配128kB内存则可支持32个虚拟机}，外加固定占用48kB程序存储器，适合中高性能微控制器使用。此外，如果多个RVM虚拟机同时使用网络协议栈等组件，使用本技术甚至可节省内存，因为网络协议栈等共享组件可以只在系统中创建一个副本，而且多个操作系统的底层切换代码和管理逻辑由于完全被抽象出去，因此在系统中只有一个副本。

在此我们回顾一下虚拟化的种类。在本手册中，我们把虚拟机按照其运行平台分成三类，分别称为原生型（I型）、宿主型（II型）和容器型（III型）。同时，依照虚拟机支持虚拟化方式的不同，又可以分为全虚拟化型（Full
Virtualization，FV）和准虚拟化型（Para-Virtualization，PV）。这两种分类方法是互相独立的。

\subsection{原生型虚拟机监视器}\label{subsec:1-2-1}

原生型虚拟机监视器作为底层软件直接运行在硬件上，创造出多个物理机的实例，也叫I型虚拟机监视器。多个虚拟机在这些物理机的实例上运行。这类虚拟机监视器的优势在于高性能、低开销，但缺点则是自身不具备虚拟机管理功能，需要0域\footnote{也即第一个启动的、具有管理特权的虚拟机}对其他虚拟机进行管理，且无法像常规操作系统那样运行轻量级的普通进程。此类虚拟机监视器有KVM\footnote{KVM模糊了原生型和宿主型的边界，严格地讲不属于其中任何一类}、Hyper-V、VMware ESXi Server、Xen、XtratuM等。

\subsection{宿主型虚拟机监视器}\label{subsec:1-2-2}

宿主型虚拟机监视器作为应用程序运行在操作系统上，创造出多个物理机的实例，也叫II型虚拟机监视器。多个虚拟机在这些物理机的实例上运行。这类虚拟机的优势在于系统独立、自身具备一定的管理功能，并允许将它们的虚拟机进程和其它普通的宿主机进程一起调度，因此得到了最广泛的应用，但效率通常认为比原生型要低一些。此类虚拟机监视器有VMware Workstation、Virtual Box、QEMU\footnote{不包括其KVM版本}等。RVM是在RME上开发的一个应用程序，其虚拟机也作为进程运行在RME上，因此也属于宿主型虚拟机监视器。

\subsection{容器型虚拟机监视器}\label{subsec:1-2-3}

容器型虚拟机监视器是由操作系统本身提供的功能，创造出多个相互隔离的操作系统资源分配空间的独立实例，严格来讲并不属于虚拟机监视器，但由于它也具有一定的资源分区功能而在某些场合被称为III型虚拟机监视器。多个应用程序直接在这些资源分配空间实例上运行。通常而言，这类虚拟机监视器的性能最高，虚拟化的内存和时间开销最低，但是它要求应用程序必须是针对它提供的接口编写的。通常而言，这接口与操作系统自身的系统调用一致，而其它接口均不被兼容，因而普适性很差。此类虚拟机监视器有Docker、Pouch、RKT、LXC等。当RVM导致的性能开销被认为不合适时，由于RME操作系统自身也原生具备容器型虚拟机监视器的能力，也可以不使用RVM而把其提供的用户态库用来开发原生的容器型虚拟机应用。事实上，如果不在RVM配置文件中声明任何虚拟机，RVM的虚拟化功能会在编译时被从系统中自动去掉以适应资源极端的环境\footnote{如Cortex-M0+等RAM在20KiB上下的环境}。

\subsection{全虚拟化型虚拟机监视器}\label{subsec:1-2-4}

全虚拟化型虚拟机监视器是有能力创建和原裸机完全一致的硬件实例的虚拟机监视器。通常而言，这需要MMU的介入来进行地址转换，而且对特权指令要进行二进制翻译、陷阱重定向或专用硬件\footnote{如Intel VT-x、AMD AMD-V等}处理，对于I/O也需要特殊硬件功能\footnote{如Intel VT-d等}来重定向。此类虚拟机在MCU环境上难于实现，因为微控制器通常都不提供MMU，也不提供相应的虚拟化硬件进行重定向功能。此类虚拟机监视器有VMware Workstation、Virtual Box、QEMU、KVM、Hyper-V、ESX Server等。

\subsection{准虚拟化型虚拟机监视器}\label{subsec:1-2-5}

准虚拟化型虚拟机监视器是创建了和原裸机有些许差别的硬件实例的虚拟机监视器。它使用超调用（Hypercall）处理特权指令和特殊I/O，并借此避开那些复杂、不必要的全虚拟化必须提供的细节。此类虚拟机在MCU环境上能够以较高的性能实现。此类虚拟机监视器有Xen等。RVM也是一个准虚拟化型虚拟机监视器，通过超调用处理系统功能。

\section{手册总览}\label{sec:1-3}

本手册从用户的角度解释了RVM的基本组成与运行原理，并假定读者对虚拟化的基本概念有一般了解。在本手册中，所有的typedef类型之前的``rvm\_''前缀均被省去。

本手册的第一部分为\protect\hyperref[sec:2-1]{第2章}
。在该章节中将概略描述RVM的架构，推荐详细阅读。

本手册的第二部分为第3章到第5章。在这些章节中将分别详细描述RVM的技术细节，推荐按需查阅。其中，\protect\hyperref[sec:3-1]{第3章}
介绍原生进程的细节信息，\protect\hyperref[sec:4-1]{第4章}
介绍虚拟机的细节信息，\protect\hyperref[sec:5-1]{第5章}
介绍工程生成器的细节信息。

本手册的第三部分为\protect\hyperref[ch:6-using-rvm]{第6章}
。在该章节中将从用户角度详述RVM及其工程生成器的使用方式，推荐详细阅读。

本手册的最后部分为\protect\hyperref[sec:8-1]{第8章}
。在该章节中将详细描述如何移植RVM到新架构。如果用户不需要进行移植，无需阅读该章节。

\needspace{12\baselineskip}
\section{本章参考文献}\label{sec:1-4}

\begin{enumerate}
\def\labelenumi{{[}\arabic{enumi}{]}}
\item
  \protect\phantomsection\label{cite:1-1}
  P. Patel, M. Vanga, and B. B. Brandenburg, ``TimerShield: protecting
  high-priority tasks from low-priority timer interference,'' in
  \textit{Proceedings of the IEEE Real-Time and Embedded Technology and
  Applications Symposium (RTAS)}, 2017, pp. 3--12.
\item
  \protect\phantomsection\label{cite:1-2}
  A. Dunkels, B. Gronvall, and T. Voigt, ``Contiki---a lightweight and
  flexible operating system for tiny networked sensors,'' in
  \textit{Proceedings of the 29th Annual IEEE International Conference
  on Local Computer Networks}, 2004, pp. 455--462.
\end{enumerate}

